\chapter{Funciones y ecuaciones de segundo grado}\label{ch:g11:quad}

Las funciones de segundo grado son las más sencillas después de las
lineales y las primeras cuyas \omterm{def:g10:functions:graph}{gráficas} son de verdad curvas. Este capítulo
desarrolla toda la caja de herramientas necesaria: la \omterm{thm:g11:quad:canonical}{forma canónica}, el
\omterm{def:g11:quad:discriminant}{discriminante}, el signo de una expresión de segundo grado y la geometría
de la \omterm{def:g11:quad:parabola}{parábola}.

\section{La forma canónica}

\begin{definition}[Función de segundo grado]\label{def:g11:quad:quadratic}
Una \emph{función de segundo grado}\index{función de segundo grado} es una
\omterm{def:g10:functions:function}{función} de la forma
\[
f(x) = ax^2 + bx + c, \qquad a \neq 0,
\]
donde $a$, $b$, $c$ son \omterm{def:g10:numbers:sets}{números reales}. La expresión $ax^2 + bx + c$ se
llama también \emph{trinomio de segundo grado}.
\end{definition}

\begin{theorem}[Forma canónica]\label{thm:g11:quad:canonical}
Toda \omterm{def:g11:quad:quadratic}{función de segundo grado} se puede escribir de una única manera como
\[
f(x) = a(x - \alpha)^2 + \beta,
\qquad\text{donde } \alpha = -\frac{b}{2a} \text{ y } \beta = f(\alpha).
\]
Es la \emph{forma canónica}\index{forma canónica} de $f$.
\end{theorem}

\begin{proof}
\emph{Completamos el cuadrado}. Como $a \neq 0$, lo sacamos como factor
común de los dos primeros términos:
\[
ax^2 + bx + c = a\left(x^2 + \frac{b}{a}x\right) + c .
\]
Ahora bien, $x^2 + \frac{b}{a}x$ es el comienzo del cuadrado
$\left(x + \frac{b}{2a}\right)^2 = x^2 + \frac ba x + \frac{b^2}{4a^2}$,
de modo que $x^2 + \frac{b}{a}x = \left(x + \frac{b}{2a}\right)^2 -
\frac{b^2}{4a^2}$. Sustituyendo:
\[
f(x) = a\left(x + \frac{b}{2a}\right)^2 - \frac{b^2}{4a} + c
= a\bigl(x - \alpha\bigr)^2 + \beta,
\]
con $\alpha = -\frac{b}{2a}$ y $\beta = c - \frac{b^2}{4a}$. Al evaluar en
$x = \alpha$ desaparece el término al cuadrado, luego $\beta = f(\alpha)$.
Unicidad: al \omterm{def:g10:algebra:expand}{desarrollar} $a(x-\alpha)^2 + \beta$ e identificar los
coeficientes de $x$ queda forzado $\alpha = -\frac{b}{2a}$, y entonces
$\beta$ está determinado.
\end{proof}

\begin{example}\label{ex:g11:quad:canonical}
Sea $f(x) = 2x^2 - 8x + 3$. Entonces $\alpha = -\frac{-8}{2 \times 2} = 2$
y $\beta = f(2) = 8 - 16 + 3 = -5$, luego
\[
f(x) = 2(x - 2)^2 - 5 .
\]
Comprobación desarrollando: $2(x^2 - 4x + 4) - 5 = 2x^2 - 8x + 8 - 5 =
2x^2 - 8x + 3$.
\end{example}

\begin{proposition}[Extremo y variación]\label{prop:g11:quad:variations}
Sea $f(x) = a(x-\alpha)^2 + \beta$.
\begin{enumerate}
  \item Si $a > 0$, entonces $f$ es \omterm{def:g10:functions:variations}{decreciente} en $\intoc{-\infty}{\alpha}$
        y \omterm{def:g10:functions:variations}{creciente} en $\intco{\alpha}{+\infty}$: $f$ tiene un
        \emph{\omterm{def:g10:functions:extrema}{mínimo}} $\beta$, alcanzado en $x = \alpha$.
  \item Si $a < 0$, la variación se invierte: $f$ tiene un \emph{\omterm{def:g10:functions:extrema}{máximo}}
        $\beta$ en $x = \alpha$.
\end{enumerate}
\end{proposition}

\begin{proof}
Supongamos $a > 0$ y sea $\alpha \leq u < v$. Entonces
\[
f(v) - f(u) = a\bigl((v-\alpha)^2 - (u-\alpha)^2\bigr)
= a\,(v - u)\,\bigl((v-\alpha) + (u-\alpha)\bigr).
\]
Cada factor es positivo: $a > 0$, $v - u > 0$ y
$(v-\alpha) + (u-\alpha) > 0$, porque $v > \alpha$ y $u \geq \alpha$.
Luego $f(v) > f(u)$: $f$ es \omterm{def:g10:functions:variations}{creciente} en $\intco{\alpha}{+\infty}$. Para
$u < v \leq \alpha$ el último factor es negativo, luego $f(v) < f(u)$: $f$
decrece ahí. Por último, $f(x) - f(\alpha) = a(x-\alpha)^2 \geq 0$ para
todo $x$, así que $\beta = f(\alpha)$ es el \omterm{def:g10:functions:extrema}{mínimo}. El caso $a < 0$ se
obtiene aplicando lo anterior a $-f$.
\end{proof}

\begin{example}
Un granjero dispone de $100$ metros de valla para cercar un campo
rectangular apoyado en un muro recto (a lo largo del muro no hace falta
valla). Si $x$ es la anchura, el área cercada es
$A(x) = x(100 - 2x) = -2x^2 + 100x$. Aquí $\alpha =
-\frac{100}{2\times(-2)} = 25$ y $A(25) = 25 \times 50 = 1250$: el área es
máxima con una anchura de $25$ m, y vale $1250$ m$^2$.
\end{example}

\section{Raíces y factorización}

\begin{definition}[Raíz, discriminante]\label{def:g11:quad:discriminant}
Una \emph{raíz}\index{raíz} del trinomio $ax^2+bx+c$ es una solución de la
\omterm{def:g10:algebra:equation}{ecuación} $ax^2 + bx + c = 0$. El
\emph{discriminante}\index{discriminante} del trinomio es el número
\[
\Delta = b^2 - 4ac .
\]
\end{definition}

\begin{theorem}[Resolución de $ax^2+bx+c=0$]\label{thm:g11:quad:roots}
\begin{enumerate}
  \item Si $\Delta > 0$, la \omterm{def:g10:algebra:equation}{ecuación} tiene dos raíces distintas
        \[
        x_1 = \frac{-b - \sqrt{\Delta}}{2a},
        \qquad
        x_2 = \frac{-b + \sqrt{\Delta}}{2a},
        \]
        y $ax^2+bx+c = a(x - x_1)(x - x_2)$.
  \item Si $\Delta = 0$, la \omterm{def:g10:algebra:equation}{ecuación} tiene la única \omterm{def:g11:quad:discriminant}{raíz}
        $x_0 = -\frac{b}{2a}$, y $ax^2+bx+c = a(x - x_0)^2$.
  \item Si $\Delta < 0$, la \omterm{def:g10:algebra:equation}{ecuación} no tiene solución real y el trinomio
        no se puede \omterm{def:g10:algebra:expand}{factorizar} con coeficientes reales.
\end{enumerate}
\end{theorem}

\begin{proof}
Partimos de la \omterm{thm:g11:quad:canonical}{forma canónica} calculada en el
\cref{thm:g11:quad:canonical}:
\[
ax^2+bx+c
= a\left(x + \frac{b}{2a}\right)^2 + c - \frac{b^2}{4a}
= a\left[\left(x + \frac{b}{2a}\right)^2 - \frac{\Delta}{4a^2}\right],
\]
ya que $c - \frac{b^2}{4a} = \frac{4ac - b^2}{4a} = -\frac{\Delta}{4a} =
a \times \left(-\frac{\Delta}{4a^2}\right)$.

\emph{1.} Si $\Delta > 0$, entonces $\frac{\Delta}{4a^2} =
\left(\frac{\sqrt\Delta}{2a}\right)^2$ y el corchete es una diferencia de
dos cuadrados:
\[
ax^2+bx+c = a\left(x + \frac{b}{2a} - \frac{\sqrt\Delta}{2a}\right)
\left(x + \frac{b}{2a} + \frac{\sqrt\Delta}{2a}\right)
= a(x - x_2)(x - x_1).
\]
Un producto es cero exactamente cuando uno de sus factores lo es, lo que
da las dos raíces, distintas porque $\sqrt\Delta \neq 0$.

\emph{2.} Si $\Delta = 0$, el corchete es
$\left(x + \frac{b}{2a}\right)^2$, que solo se anula en
$x_0 = -\frac{b}{2a}$.

\emph{3.} Si $\Delta < 0$, entonces $-\frac{\Delta}{4a^2} > 0$, así que el
corchete es un cuadrado más un número positivo: no se anula nunca y la
\omterm{def:g10:algebra:equation}{ecuación} no tiene solución real. Una factorización $a(x-x_1)(x-x_2)$
produciría raíces, luego no existe ninguna.
\end{proof}

\begin{example}\label{ex:g11:quad:solving}
Resolvamos $2x^2 - 8x + 3 = 0$. Aquí
$\Delta = (-8)^2 - 4 \times 2 \times 3 = 64 - 24 = 40 > 0$, así que hay
dos raíces:
\[
x_{1,2} = \frac{8 \pm \sqrt{40}}{4} = \frac{8 \pm 2\sqrt{10}}{4}
= 2 \pm \frac{\sqrt{10}}{2}.
\]
Para $x^2 + x + 1 = 0$: $\Delta = 1 - 4 = -3 < 0$, sin solución real.
\end{example}

\begin{proposition}[Suma y producto de las raíces]\label{prop:g11:quad:vieta}
Si $\Delta \geq 0$ y $x_1, x_2$ son las raíces (iguales cuando
$\Delta = 0$), entonces
\[
x_1 + x_2 = -\frac{b}{a},
\qquad
x_1 x_2 = \frac{c}{a} .
\]
Recíprocamente, dos números de suma $s$ y producto $p$ son las soluciones
de $x^2 - sx + p = 0$.
\end{proposition}

\begin{proof}
Desarrollamos la factorización del \cref{thm:g11:quad:roots}:
\[
a(x - x_1)(x - x_2) = a x^2 - a(x_1 + x_2)x + a\,x_1 x_2 .
\]
Identificando los coeficientes con los de $ax^2 + bx + c$ se obtiene
$-a(x_1+x_2) = b$ y $a\,x_1x_2 = c$. Para el recíproco: si $u + v = s$ y
$uv = p$, entonces $(x-u)(x-v) = x^2 - sx + p$, así que $u$ y $v$ son sus
raíces.
\end{proof}

\begin{example}
Hallemos dos números de suma $10$ y producto $21$. Son las soluciones de
$x^2 - 10x + 21 = 0$; $\Delta = 100 - 84 = 16$ y las raíces son
$\frac{10 \pm 4}{2} = 7$ y $3$. \emph{Raíces a simple vista:} como los
coeficientes de $x^2 - 5x + 4$ cumplen $1 - 5 + 4 = 0$, el número $1$ es
\omterm{def:g11:quad:discriminant}{raíz}, y la otra es $\frac{c}{a} = 4$ (su producto debe valer
$\frac ca$).
\end{example}

\section{Signo de un trinomio de segundo grado}

\begin{theorem}[Signo de un trinomio]\label{thm:g11:quad:sign}
Sea $f(x) = ax^2+bx+c$ con $a \neq 0$.
\begin{enumerate}
  \item Si $\Delta > 0$, con raíces $x_1 < x_2$: $f(x)$ tiene el signo de
        $a$ fuera de $\intcc{x_1}{x_2}$ y el signo de $-a$ en
        $\intoo{x_1}{x_2}$.
  \item Si $\Delta = 0$: $f(x)$ tiene el signo de $a$ para todo
        $x \neq x_0$, y se anula en $x_0$.
  \item Si $\Delta < 0$: $f(x)$ tiene el signo de $a$ para todo $x$.
\end{enumerate}
En resumen: \emph{un trinomio tiene el signo de $a$, salvo entre sus
raíces}.
\end{theorem}

\begin{proof}
\emph{1.} Escribimos $f(x) = a(x-x_1)(x-x_2)$ y seguimos los signos de los
factores. Si $x < x_1$: tanto $x - x_1$ como $x - x_2$ son negativos, su
producto es positivo y $f(x)$ tiene el signo de $a$. Si $x_1 < x < x_2$:
los factores tienen signos opuestos, el producto es negativo y $f(x)$
tiene el signo de $-a$. Si $x > x_2$: los dos factores son positivos y
$f(x)$ vuelve a tener el signo de $a$.

\emph{2.} $f(x) = a(x - x_0)^2$ y $(x-x_0)^2 > 0$ para $x \neq x_0$.

\emph{3.} Por la demostración del \cref{thm:g11:quad:roots},
$f(x) = a\left[\left(x + \frac{b}{2a}\right)^2 +
\left(-\frac{\Delta}{4a^2}\right)\right]$, y el corchete es siempre
positivo.
\end{proof}

\begin{method}[Resolver inecuaciones de segundo grado]\label{met:g11:quad:inequalities}
Para resolver $ax^2 + bx + c \leq 0$ (o $\geq$, $<$, $>$):
\begin{enumerate}
  \item calcula $\Delta$ y las raíces, si las hay;
  \item construye la \omterm{def:g10:algebra:signtable}{tabla de signos} usando el \cref{thm:g11:quad:sign}
        (signo de $a$ fuera de las raíces);
  \item lee el \omterm{def:g10:numbers:interval}{intervalo} o los \omterm{def:g10:numbers:interval}{intervalos} solución, atento a si las raíces
        están incluidas (desigualdad no estricta) o no (estricta).
\end{enumerate}
\end{method}

\begin{example}\label{ex:g11:quad:inequality}
Resolvamos $-x^2 + 3x + 4 \geq 0$. Aquí $\Delta = 9 + 16 = 25$ y las
raíces son $\frac{-3 \pm 5}{-2}$, es decir, $x_1 = -1$ y $x_2 = 4$. Como
$a = -1 < 0$, el trinomio es negativo fuera de las raíces y positivo entre
ellas. Con la desigualdad no estricta, el conjunto de soluciones es
$\intcc{-1}{4}$.
\end{example}

\section{La parábola}

\begin{definition}[Parábola]\label{def:g11:quad:parabola}
La \omterm{def:g10:functions:graph}{gráfica} de una \omterm{def:g11:quad:quadratic}{función de segundo grado} $f(x) = a(x-\alpha)^2+\beta$ es
una \emph{parábola}\index{parábola} de \emph{vértice} $S(\alpha, \beta)$.
Se abre hacia arriba si $a > 0$ y hacia abajo si $a < 0$, y es simétrica
respecto de la recta vertical $x = \alpha$.
\end{definition}

\begin{remark}
La simetría es fácil de comprobar: para cualquier $h$,
$f(\alpha + h) = ah^2 + \beta = f(\alpha - h)$, así que dos puntos a la
misma distancia horizontal de la recta $x = \alpha$ tienen la misma
altura.
\end{remark}

\begin{omfigure}
\begin{tikzpicture}
\begin{axis}[omaxis,
  width=5.1cm, height=4.8cm,
  xmin=-0.8, xmax=4.4, ymin=-1.7, ymax=3.6,
  xtick={1,3}, xticklabels={$x_1$,$x_2$}, ytick=\empty,
  title={$\Delta > 0$}]
  \addplot[omDef, thick, domain=-0.15:4.15] {(x-1)*(x-3)};
  \fill (axis cs:1,0) circle (1.5pt);
  \fill (axis cs:3,0) circle (1.5pt);
  \fill (axis cs:2,-1) circle (1.5pt) node[below, font=\small] {$S$};
\end{axis}
\end{tikzpicture}%
\begin{tikzpicture}
\begin{axis}[omaxis,
  width=5.1cm, height=4.8cm,
  xmin=-0.8, xmax=4.4, ymin=-1.7, ymax=3.6,
  xtick={2}, xticklabels={$x_0$}, ytick=\empty,
  title={$\Delta = 0$}]
  \addplot[omDef, thick, domain=0.2:3.8] {(x-2)^2};
  \fill (axis cs:2,0) circle (1.5pt)
    node[above right=1pt, font=\small] {$S$};
\end{axis}
\end{tikzpicture}%
\begin{tikzpicture}
\begin{axis}[omaxis,
  width=5.1cm, height=4.8cm,
  xmin=-0.8, xmax=4.4, ymin=-1.7, ymax=3.6,
  xtick=\empty, ytick=\empty,
  title={$\Delta < 0$}]
  \addplot[omDef, thick, domain=0.35:3.65] {(x-2)^2 + 0.8};
  \fill (axis cs:2,0.8) circle (1.5pt) node[below, font=\small] {$S$};
\end{axis}
\end{tikzpicture}

{\small Las tres posiciones de una \omterm{def:g11:quad:parabola}{parábola} abierta hacia arriba
($a > 0$) respecto del eje $x$: dos raíces, una \omterm{def:g11:quad:discriminant}{raíz} doble, ninguna \omterm{def:g11:quad:discriminant}{raíz}.
El vértice $S$ tiene \omterm{def:g10:coordgeom:system}{abscisa} $\alpha = -\frac{b}{2a}$.}
\end{omfigure}

\begin{omfigure}
\begin{tikzpicture}
\begin{axis}[omaxis,
  width=8.6cm, height=5.6cm,
  xmin=-2.2, xmax=5.4, ymin=-4.4, ymax=6.2,
  xtick={-1,4}, xticklabels={$x_1$,$x_2$}, ytick=\empty]
  \addplot[name path=par, omDef, thick, domain=-1.75:4.75, samples=80]
    {-0.8*(x+1)*(x-4)};
  \path[name path=xaxis] (-1.75,0) -- (4.75,0);
  \addplot[omDef!20] fill between[of=par and xaxis,
    soft clip={domain=-1:4}];
  \fill (axis cs:-1,0) circle (1.5pt);
  \fill (axis cs:4,0) circle (1.5pt);
  \node[font=\small, omDef] at (axis cs:1.5,2.0) {$f > 0$};
  \node[font=\small, omThm] at (axis cs:-1.75,-1.3) {$f < 0$};
  \node[font=\small, omThm] at (axis cs:4.75,-1.3) {$f < 0$};
\end{axis}
\end{tikzpicture}

{\small Signo de un trinomio abierto hacia abajo ($a < 0$) con dos raíces:
el signo de $-a$ entre las raíces (zona sombreada) y el de $a$ fuera.}
\end{omfigure}

\begin{method}[Leer una parábola]\label{met:g11:quad:reading}
Para esbozar $f(x) = ax^2+bx+c$: halla la orientación (signo de $a$), el
vértice $\left(-\frac{b}{2a},\, f\!\left(-\frac{b}{2a}\right)\right)$, el
corte $(0, c)$ con el eje $y$ y las raíces, si las hay. El eje de simetría
$x = -\frac{b}{2a}$ coloca después dos veces cada punto calculado.
\end{method}

\begin{example}
Para $f(x) = x^2 - 4x + 3$: abierta hacia arriba, vértice $(2, -1)$, corta
al eje $y$ en $(0,3)$ y tiene raíces $1$ y $3$ (a simple vista:
$1 - 4 + 3 = 0$). Por simetría, el punto $(4, 3)$ también está en la
curva, reflejando a $(0, 3)$.
\end{example}

\section{Ejercicios}

\begin{exercise}[$\star$]\label{exo:g11:quad:1}
Resuelve las ecuaciones
\[
x^2 - 5x + 6 = 0, \qquad
2x^2 + 3x - 2 = 0, \qquad
x^2 + 2x + 5 = 0, \qquad
9x^2 - 6x + 1 = 0 .
\]
\end{exercise}

\begin{exercise}[$\star$]\label{exo:g11:quad:2}
Escribe en \omterm{thm:g11:quad:canonical}{forma canónica} y da el vértice de la \omterm{def:g11:quad:parabola}{parábola}:
\[
f(x) = x^2 - 6x + 11, \qquad
g(x) = -2x^2 + 4x + 1, \qquad
h(x) = 3x^2 + 6x .
\]
\end{exercise}

\begin{exercise}[$\star$]\label{exo:g11:quad:3}
Resuelve las inecuaciones
\[
x^2 - 3x - 10 < 0, \qquad
2x^2 + x + 3 > 0, \qquad
-x^2 + 6x - 9 \geq 0 .
\]
\end{exercise}

\begin{exercise}[$\star$]\label{exo:g11:quad:4}
Sin calcular $\Delta$, halla una \omterm{def:g11:quad:discriminant}{raíz} evidente y después la otra:
\[
x^2 - 7x + 6 = 0, \qquad
2x^2 + 5x - 7 = 0, \qquad
x^2 + 4x - 5 = 0 .
\]
(Indicación: prueba $x = 1$ y $x = -1$ y usa después el producto de las
raíces.)
\end{exercise}

\begin{exercise}[$\star\star$]\label{exo:g11:quad:5}
Halla dos números cuya suma sea $14$ y cuyo producto sea $45$. Halla
después las dimensiones de un rectángulo de perímetro $34$ y área $60$.
\end{exercise}

\begin{exercise}[$\star\star$]\label{exo:g11:quad:6}
Sea $f(x) = x^2 + (m-2)x + 1$, donde $m$ es un parámetro real. ¿Para qué
valores de $m$ tiene la \omterm{def:g10:algebra:equation}{ecuación} $f(x) = 0$ exactamente una solución? ¿Y
ninguna?
\end{exercise}

\begin{exercise}[$\star\star$]\label{exo:g11:quad:7}
Se lanza una pelota hacia arriba; su altura al cabo de $t$ segundos es
$h(t) = -5t^2 + 20t + 1$ (en metros). ¿Cuál es la altura máxima alcanzada
y en qué instante? ¿Durante qué \omterm{def:g10:numbers:interval}{intervalo} de tiempo está la pelota al
menos a $16$ m de altura?
\end{exercise}

\begin{exercise}[$\star\star$]\label{exo:g11:quad:8}
Resuelve la \omterm{def:g10:algebra:equation}{ecuación} $x^4 - 13x^2 + 36 = 0$ tomando $X = x^2$ (una
\omterm{def:g10:algebra:equation}{ecuación} \emph{bicuadrada}).
\end{exercise}

\begin{exercise}[$\star\star$]\label{exo:g11:quad:9}
¿Para qué valores de $x$ está el punto $M(x, x^2)$ de la \omterm{def:g11:quad:parabola}{parábola}
$y = x^2$ estrictamente por debajo de la recta $y = x + 2$? Interprétalo
sobre un esbozo.
\end{exercise}

\begin{exercise}[$\star\star$]\label{exo:g11:quad:10}
La suma de un número positivo y su inverso es $\frac{13}{6}$. Halla el
número. (Redúcelo a una \omterm{def:g10:algebra:equation}{ecuación} de segundo grado.)
\end{exercise}

\begin{exercise}[$\star\star\star$]\label{exo:g11:quad:11}
Sea $\mathcal P$ la \omterm{def:g11:quad:parabola}{parábola} $y = x^2$ y sea $d_m$ la recta $y = mx - 1$,
donde $m$ es un parámetro real.
\begin{enumerate}
  \item ¿Para qué valores de $m$ se cortan $\mathcal P$ y $d_m$ en dos
        puntos? ¿En un punto? ¿En ninguno?
  \item Cuando se tocan exactamente en un punto, calcula el punto de
        contacto. (En el \cref{ch:g11:deriv} reconocerás $d_m$ como la
        tangente a $\mathcal P$ en ese punto.)
\end{enumerate}
\end{exercise}

\section{Problema: El punto secreto de la parábola}

\begin{problem}[{Problema de fin de semana --- foco y directriz: por qué
las antenas parabólicas, los faros de los coches y los puentes colgantes
dibujan todos la misma curva, de la que en el fondo solo hay una}]
\label{pb:g11:quad:1}
Todas las antenas parabólicas de todos los balcones apuntan su receptor a
un punto interior muy concreto; todos los faros de coche esconden su
bombilla en ese mismo lugar especial. Ese punto, el \emph{\omterm{pb:g11:quad:1}{foco}}, es un
secreto que la \omterm{def:g11:quad:parabola}{parábola} guarda desde la geometría griega, y la fórmula de
la distancia basta para arrancárselo. De camino, este problema entrena el
\omterm{def:g11:quad:discriminant}{discriminante}, lanza proyectiles por encima de muros, cuelga un puente y
demuestra un hecho final asombroso: salvo aumento, \emph{en el mundo solo
hay una \omterm{def:g11:quad:parabola}{parábola}}.

\medskip
\noindent\textbf{Parte I --- Soltura con los trinomios.}
\begin{enumerate}
  \item Resuelve $x^2 - 5x + 6 = 0$, después $2x^2 + x - 6 = 0$ y después
        $x^2 + x + 1 = 0$ (\cref{thm:g11:quad:roots}).
  \item Escribe $2x^2 - 4x - 6$ en forma factorizada y en \omterm{thm:g11:quad:canonical}{forma canónica}
        (\cref{thm:g11:quad:canonical}). Junto con la forma desarrollada,
        tienes ya tres disfraces de una misma \omterm{def:g10:functions:function}{función}: di qué rasgo enseña
        cada uno de un vistazo (raíces, vértice, corte con el eje $y$).
  \item Resuelve $-x^2 + 4x - 3 \geq 0$ con la regla de los signos
        (\cref{thm:g11:quad:sign}).
  \item Viète como detective (\cref{prop:g11:quad:vieta}, continuando
        Babilonia en el \cref{pb:g10:algebra:1}): halla dos números de
        suma $7$ y producto $12$ escribiendo su \omterm{def:g10:algebra:equation}{ecuación}. Después: una
        \omterm{def:g11:quad:discriminant}{raíz} de $x^2 - 5x + c = 0$ vale $2$; halla $c$ y la otra \omterm{def:g11:quad:discriminant}{raíz} sin
        resolver nada.
  \item ¿Para qué valores de $m$ tiene $x^2 + mx + 9 = 0$ una \omterm{def:g11:quad:discriminant}{raíz} doble?
        ¿Ninguna \omterm{def:g11:quad:discriminant}{raíz}? ¿Dos raíces?
\end{enumerate}

\medskip
\noindent\textbf{Parte II --- El \omterm{pb:g11:quad:1}{foco} al descubierto.}
Trabajamos con la \omterm{def:g11:quad:parabola}{parábola} $y = x^2$, el punto
$F\left(0, \frac14\right)$ y la recta horizontal $\delta$ de \omterm{def:g10:algebra:equation}{ecuación}
$y = -\frac14$.
\begin{enumerate}[resume]
  \item Para un punto $P(x, x^2)$ de la \omterm{def:g11:quad:parabola}{parábola}, desarrolla
        $PF^2 = x^2 + \left(x^2 - \frac14\right)^2$ y demuestra que es el
        cuadrado perfecto $\left(x^2 + \frac14\right)^2$. Deduce $PF$.
  \item Calcula la distancia de $P$ a la recta $\delta$ y concluye:
        \emph{todo punto de la \omterm{def:g11:quad:parabola}{parábola} está exactamente a la misma
        distancia de $F$ que de $\delta$}. ($F$ es el
        \emph{foco}\index{foco} y $\delta$, la
        \emph{directriz}\index{directriz}.)
  \item Demuestra el recíproco: un punto $P(x, y)$ con
        $PF = \operatorname{dist}(P, \delta)$ tiene que cumplir
        $y = x^2$. (Eleva al cuadrado las dos distancias y simplifica.)
        Así pues, la \omterm{def:g11:quad:parabola}{parábola} es un \emph{lugar geométrico}, exactamente
        como lo era la circunferencia en el \cref{pb:g10:coordgeom:1}:
        equidistancia de un punto y una recta, en lugar de equidistancia
        de un solo punto.
  \item Para la \omterm{def:g11:quad:parabola}{parábola} más ancha $y = \frac{x^2}{4f}$ (con $f > 0$),
        adapta la pregunta 6 para demostrar que el \omterm{pb:g11:quad:1}{foco} es $(0, f)$ y la
        \omterm{pb:g11:quad:1}{directriz}, $y = -f$. Aplicación: una antena parabólica mide $60$
        cm de ancho y $9$ cm de profundidad, así que su borde pasa por
        $(30, 9)$. Halla $f$: ¿a qué altura sobre el vértice hay que
        colgar el receptor?
  \item La razón de que las antenas funcionen: los rayos que llegan
        paralelos al eje se reflejan todos \emph{pasando por el \omterm{pb:g11:quad:1}{foco}} (y
        una bombilla en el \omterm{pb:g11:quad:1}{foco} lanza un haz paralelo: los faros son
        antenas al revés). Explica por qué la equidistancia de la pregunta
        7 hace esto verosímil (piensa en los rayos como frentes de onda
        que avanzan) y di qué herramienta, que llega en el
        \cref{ch:g11:deriv} (véase el \cref{exo:g11:quad:11}), convierte
        la verosimilitud en demostración.
\end{enumerate}

\medskip
\noindent\textbf{Parte III --- \omterm{def:g11:quad:parabola}{Parábolas} trabajando.}
\begin{enumerate}[resume]
  \item Una pelota sigue la trayectoria $y = x - \frac{x^2}{20}$ ($x$ e
        $y$ en metros). Halla dónde cae, la posición horizontal de su
        punto más alto y la altura máxima
        (\cref{prop:g11:quad:variations}).
  \item En la misma trayectoria hay un muro en $x = 4$. Calcula la altura
        de la pelota allí: ¿supera un muro de $3$ m? ¿Y uno de $4$ m?
  \item Un arco parabólico salva $40$ m con una altura máxima de $10$ m:
        en \omterm{def:g10:coordgeom:system}{coordenadas} centradas, $y = 10 - \frac{x^2}{40}$. ¿Puede pasar
        un camión de $7$ m de alto a $10$ m del centro? ¿Y a $15$ m?
  \item Los ingresos de un teatro con un precio de entrada de $p$ euros
        son $R(p) = p\,(120 - 2p)$ (la demanda baja al subir el precio).
        Halla el precio que maximiza los ingresos y el \omterm{def:g10:functions:extrema}{máximo}.
  \item Los cables principales del Golden Gate son parabólicos con
        excelente \omterm{def:g10:numbers:approx}{aproximación} (una carga uniforme del tablero produce una
        \omterm{def:g11:quad:parabola}{parábola}; una cadena que cuelga libremente, \emph{no}: su curva,
        la catenaria, necesita la exponencial del año que viene). Con una
        luz de $1\,280$ m y una flecha de $143$ m, el cable sigue
        $y = 143\left(\frac{x}{640}\right)^2$ desde el centro. ¿A qué
        altura sobre su punto más bajo está el cable en $x = 320$ m?
\end{enumerate}

\medskip
\noindent\textbf{Parte IV --- Rectas, tangentes y la última
sorpresa.}
\begin{enumerate}[resume]
  \item Completa el \cref{exo:g11:quad:11}: las rectas $y = mx - 1$ que
        pasan por el punto exterior $(0, -1)$ cortan a $y = x^2$ según el
        signo de un \omterm{def:g11:quad:discriminant}{discriminante}. Halla las dos rectas tangentes desde
        $(0, -1)$ y sus puntos de contacto.
  \item Ahora las rectas $y = mx + 1$ que pasan por el punto
        \emph{interior} $(0, 1)$: calcula el \omterm{def:g11:quad:discriminant}{discriminante} de la \omterm{def:g10:algebra:equation}{ecuación}
        de los cortes y demuestra que todas esas rectas cortan dos veces a
        la \omterm{def:g11:quad:parabola}{parábola}. Moraleja: desde dentro no hay tangentes, igual que en
        una circunferencia.
  \item La cuerda focal: corta $y = x^2$ con la recta horizontal que pasa
        por el \omterm{pb:g11:quad:1}{foco}, $y = \frac14$. ¿Cuánto mide la cuerda? Comprueba la
        regla general (que conviene recordar para las antenas): para
        $y = \frac{x^2}{4f}$, la cuerda que pasa por el \omterm{pb:g11:quad:1}{foco} paralela a la
        \omterm{pb:g11:quad:1}{directriz} mide $4f$.
  \item La última sorpresa: aplica el aumento
        $(x, y) \mapsto (kx, ky)$ a la \omterm{def:g11:quad:parabola}{parábola} $y = x^2$ y demuestra que
        la \omterm{def:g10:functions:function}{imagen} es exactamente $y = \frac{x^2}{k}$. Concluye: toda
        \omterm{def:g11:quad:parabola}{parábola} $y = \frac{x^2}{4f}$ es una ampliación de la \omterm{def:g11:quad:parabola}{parábola}
        unidad; a diferencia de las circunferencias y de las elipses de
        proporciones distintas, \emph{todas las \omterm{def:g11:quad:parabola}{parábolas} son
        semejantes}.
  \item Final: el retrato de la \omterm{def:g11:quad:parabola}{parábola} en cinco líneas: una \omterm{def:g10:algebra:equation}{ecuación} con
        tres disfraces (pregunta 2), un lugar geométrico (preguntas 7 y
        8), el reflector del ingeniero (preguntas 9 y 10), la trayectoria
        de un proyectil y el cable de un puente (preguntas 11 a 15) y una
        única curva salvo aumento (pregunta 19). Una frase para cada una.
\end{enumerate}
\end{problem}
